\documentclass[11pt]{article}
\usepackage[utf8]{inputenc}
\usepackage[T1]{fontenc}
\usepackage[italian]{babel}
\usepackage{geometry}
\geometry{a4paper, margin=1in}
\usepackage{xcolor}
\usepackage{listings}
\usepackage{hyperref}

\usepackage{tcolorbox}
\tcbuselibrary{breakable}
\begin{document}

\begin{tcolorbox}[colback=blue!5!white,colframe=blue!75!black,title=\textbf{DOMANDA}, breakable]
13) Provare a rifare il caso dell'Esercizio 10 ma con questa nuova versione della sommatrice.
Cosa si può osservare? Che soluzione si può trovare? C'è influenza reciproca tra i due client?
\end{tcolorbox}

\begin{tcolorbox}[colback=green!5!white,colframe=green!75!black,title=\textbf{RISPOSTA}, breakable]
Rifacendo il caso dell'Esercizio 10 con la nuova versione \texttt{TCP} della sommatrice, ci troveremo di fronte a uno scenario completamente diverso rispetto a \texttt{UDP}.

Ecco cosa succede rispondendo alle tue domande:

\begin{itemize}
    \item \textbf{Cosa si può osservare?} Osserverai che il server gestirà soltanto il Client A, ignorando del tutto il Client B fino alla fine. Quando avvii il server, questo si blocca sulla \texttt{acceptConnection()} in attesa del primo client. Avviando il Client A, la connessione viene instaurata e il server entra nel ciclo per calcolare la somma di A. Se nel frattempo avvii il Client B, a livello di rete il sistema operativo del server accetterà "dietro le quinte" la connessione \texttt{TCP} (grazie alla coda di \texttt{listen()}), permettendo al Client B di avviarsi e farti inserire i numeri. Tuttavia, il programma del server è bloccato all'interno del ciclo \texttt{do-while} a leggere esclusivamente dal canale (socket) del Client A. I numeri inviati dal Client B si accumuleranno semplicemente nel buffer del sistema operativo in attesa di essere letti. Quando il Client A invierà lo "0", il server gli risponderà con la somma corretta (es. 7535), chiuderà la connessione e, per come è scritto attualmente il codice, il programma del server terminerà. Di conseguenza, il Client B non riceverà mai alcuna risposta, la sua connessione verrà interrotta bruscamente dal sistema operativo (generando un errore se prova a inviare o ricevere altro) e i suoi dati andranno persi senza essere stati elaborati.
    
    \item \textbf{C'è influenza reciproca tra i due client?} Rispetto ai dati, assolutamente no. A differenza di \texttt{UDP}, dove i numeri dei due client finivano mischiati nello stesso "calderone", \texttt{TCP} crea un "tubo" dedicato per ogni connessione. I dati non si sovrappongono in alcun modo e le somme non si inquinano a vicenda. L'unica (pesante) influenza reciproca è di tipo temporale/bloccante: il Client A "monopolizza" il server impedendo che il Client B venga servito (e nel nostro caso, causandone poi l'esclusione definitiva).
    
    \item \textbf{Che soluzione si può trovare?} Per risolvere il problema ci sono due step:
    \begin{itemize}
        \item \textbf{Per servirli in modo sequenziale (uno dopo l'altro):} È necessario racchiudere la \texttt{acceptConnection} e la logica della somma all'interno di un ciclo \texttt{while(true)} nel server. In questo modo, dopo aver servito il Client A (alla ricezione dello "0"), il server tornerà in attesa e farà la \texttt{accept()} della connessione del Client B che era rimasta "parcheggiata" in coda.
        \item \textbf{Per servirli in modo concorrente (contemporaneamente):} Serve un'architettura multi-thread (o multi-processo). Il server accetta la connessione nel thread principale e poi passa il socket appena creato a un nuovo thread "figlio" (es. usando la \texttt{createConcurrentNetworkFunction} fornita dalla libreria \texttt{network.c}). Il thread figlio si occuperà di fare la somma per quel client, mentre il padre tornerà immediatamente in attesa di nuove connessioni (\texttt{acceptConnection}) per altri client.
    \end{itemize}
\end{itemize}
\end{tcolorbox}

\end{document}
